
\subsubsection{4.1 Research Questions}

\begin{itemize}
\item \textbf{RQ1 (Mechanism):} Does error type determine localization reliability, and does unreliable localization cause gradient noise?
\item \textbf{RQ2 (Efficiency):} Does error-type gating improve training efficiency?
\end{itemize}

\subsubsection{4.2 Hypothesis Structure}

We decomposed RQ1 into four mechanism hypotheses:

\begin{table}[htbp]
\centering
\resizebox{\columnwidth}{!}{%
\begin{tabular}{lll}
\toprule
ID & Hypothesis & Gate Type \\
\midrule
H-M1 & Fine-grained penalties concentrate gradients at traceback locations & MUST\_WORK \\
H-M2 & Localization accuracy varies by error type & SHOULD\_WORK \\
H-M3 & Unreliable localization causes gradient concentration at wrong tokens & MUST\_WORK \\
H-M4 & Gating improves gradient signal-to-noise ratio & SHOULD\_WORK \\
\bottomrule
\end{tabular}
}
\end{table}

For RQ2:

\begin{table}[htbp]
\centering
\resizebox{\columnwidth}{!}{%
\begin{tabular}{lll}
\toprule
ID & Hypothesis & Gate Type \\
\midrule
H-E1 & Error-type gating improves sample efficiency & MUST\_WORK \\
\bottomrule
\end{tabular}
}
\end{table}

\subsubsection{4.3 Dataset}

APPS dataset was used for all experiments. For mechanism experiments (H-M1 through H-M4), 500 samples per error category were used. The efficiency experiment (H-E1) used a subset for proof-of-concept validation.

\subsubsection{4.4 Model}

CodeT5-small (60M parameters) was used for proof-of-concept validation. RLTF uses CodeT5-large (770M parameters); the smaller model enables rapid iteration but results may not directly transfer to larger scales.

\subsubsection{4.5 Ground Truth Annotation}

For H-M2 and H-M3, ground-truth bug locations are required to measure localization accuracy and gradient concentration at correct tokens. We used synthetic code templates with known bug locations, ensuring deterministic annotation. This simplifies validation but limits generalization claims.

\subsubsection{4.6 Metrics}

\begin{itemize}
\item \textbf{Concentration ratio} (H-M1, H-M3): Gradient magnitude at error-line tokens divided by magnitude at non-error tokens
\item \textbf{Localization accuracy} (H-M2): Percentage of samples where traceback line matches ground-truth bug location (within $\pm 2$ lines tolerance)
\item \textbf{Signal-to-noise ratio} (H-M4): Gradient signal at ground-truth locations divided by signal at incorrect locations
\item \textbf{Steps to threshold} (H-E1): Training steps required to reach 30\% pass@1
\end{itemize}

\subsubsection{4.7 Conditions}

Three feedback conditions were compared:

\begin{itemize}
\item \textbf{Coarse-only:} Binary pass/fail reward, no token-level penalties
\item \textbf{Fine-always:} Standard RLTF with uniform fine-grained penalties
\item \textbf{Fine-gated:} Fine-grained for U\_line, coarse for U\_ignore
\end{itemize}

